%% ============================================================================
%%  LAFROX-Subdocumento1-Anexos.tex -- Anexo 1.A: Tabla de trazabilidad
%%  observación–respuesta–sección modificada de las observaciones del Informe 2
%%  (Art. 46 de las Bases Administrativas).
%%
%%  Las tablas se generan en partes/tablas.tex y partes/tablas_sintesis.tex.
%%  Compilar:  latexmk LAFROX-Subdocumento1-Anexos.tex   (LuaLaTeX, ver latexmkrc)
%% ============================================================================
\documentclass[carta,firma]{lafrox}
\usepackage{needspace}
\usepackage{xurl}

\datoslafrox{
  caso         = {Caso --- Logística},
  fecha        = {11 de noviembre de 2026},
  subtitulo    = {Subdocumento 01 --- Anexos},
  alcance      = {Subdocumentos 1 a 9 y 13},
  formulario   = {Anexo 1.A --- Tabla de trazabilidad (Art. 46)},
  nomenclatura = {LAFROX-Subdocumento1-Anexos.pdf},
}

\configurarportadalafrox{Subdocumento 01 \textemdash{} Anexos}{Anexo 1.A \textemdash{} Trazabilidad de observaciones del Informe 2}{Art. 46}{}{FIRMA_Alex_Aravena_transparente.png}
\makeatletter\def\LFXfecha{11 de noviembre de 2026}\makeatother

\AtBeginDocument{\hypersetup{pdftitle={Anexo 1.A: Trazabilidad de observaciones del Informe 2},pdfauthor={LafroX SpA}}}
\renewcommand{\thesection}{1.A.\arabic{section}}
\renewcommand{\thesubsection}{\thesection.\arabic{subsection}}
\renewcommand{\thetable}{\arabic{table}}
\setcounter{tocdepth}{2}

\begin{document}
\portadalafrox
\indicelafrox

\Needspace{6\baselineskip}
\section{Alcance y criterio de lectura}

Este anexo responde el Art. 46 de las Bases Administrativas, que exige resolver en forma explícita las observaciones de la instancia anterior mediante una tabla observación–respuesta–sección modificada (Distribuidora Puelche S.A., 2026a, Art. 46, p. 28). Reúne las observaciones de la primera revisión del Informe 2 a los Subdocumentos 1 a 9 y 13 y a sus Formularios T-6, T-9 a T-19, más las observaciones generales de formalidad. Reemplaza al archivo de trazabilidad del Informe 1, cuyo nombre no seguía la nomenclatura de la sección 2 de la Aclaración N.° 2.

Se registran las observaciones que piden un cambio: las marcadas como «Parcial», «No corregida» o «Crítico» y las que la revisión formula sin calificativo. Los aspectos que la revisión calificó como «OK» o «Corregida» no se repiten. Cuando una misma observación reaparece en «Qué se espera en la Propuesta Final» o en las consideraciones transversales, se registra una sola vez, en el subdocumento donde se corrige.

Cada respuesta comienza con su estado. «Corregida» indica que el cambio está en la versión vigente del subdocumento o formulario. «Parcial», que el cambio atiende una parte de la observación y declara lo que falta. «Pendiente», que la versión vigente aún no contiene la corrección. La columna «Sección modificada» indica el lugar donde se lee la respuesta. El identificador combina el subdocumento y un número correlativo: SD3-04 es la cuarta observación registrada del Subdocumento 3; S41 y S42 corresponden a las secciones 4.1 y 4.2–4.3 del Subdocumento 4, y GEN a las observaciones generales.

La Tabla~\ref{tab:sintesis} resume el estado por subdocumento y las Tablas 2 a 13 presentan el detalle. La mayor parte de las observaciones quedó corregida en los subdocumentos que venían del Informe 1 y en los Subdocumentos 7, 8 y 9. El Subdocumento 6 se completó con las reglas de gestión y la adaptación al caso que faltaban. Las pendientes se concentran en el Subdocumento 7 y en puntos acotados de la formalidad general y de los Subdocumentos 3, 4, 5 y 8.

\input{partes/tablas_sintesis.tex}

\clearpage
\Needspace{6\baselineskip}
\section{Observaciones, respuestas y secciones modificadas}
\input{partes/tablas.tex}

\microtypesetup{protrusion=false}
\Needspace{6\baselineskip}
\section*{Referencias}
\addcontentsline{toc}{section}{Referencias}
\begin{itemize}[itemsep=1pt]
\item Distribuidora Puelche S.A. (2026a). \emph{Bases Administrativas de Licitación N.º TFEP-01/2026: Contratación de Solución Integral de Software y Servicios de Operación}.
\item Distribuidora Puelche S.A. (2026b). \emph{Bases Técnicas Transversales de Licitación N.º TFEP-01/2026}.
\item Distribuidora Puelche S.A. (2026c). \emph{Bases Técnicas del caso: Caso 02 --- Logística}.
\end{itemize}

\Needspace{10\baselineskip}
\section*{Declaración de uso de IA}
\addcontentsline{toc}{section}{Declaración de uso de IA}

La Tabla~\ref{tab:declaracion_ia} identifica la herramienta y el alcance de la asistencia utilizada en este anexo.

\begin{tablalafrox}{Uso de IA en el Anexo 1.A. Fuente: registro del equipo.}{tab:declaracion_ia}{L{1.6cm} L{2.0cm} Y L{1.2cm} L{1.6cm} L{3.4cm}}
  {\cab{Sección} & \cab{Herramienta} & \cab{Finalidad del uso} & \cab{Nivel en texto} & \cab{Nivel en diagramas} & \cab{Revisión humana (quién y qué verificó)}}
  1.A.1 & Claude Code & Redacción del alcance y de la síntesis & Alto & Ninguno & Guillermo Castillo: criterio de registro y conteo de estados. \\
  1.A.2 & Claude Code & Cotejo de cada observación con la versión corregida de los subdocumentos y redacción de las respuestas & Alto & Ninguno & Guillermo Castillo: estado y sección modificada de cada fila contra el subdocumento correspondiente. \\
\end{tablalafrox}

\end{document}
